% Chapitre 3 — Analyse et conception
\chapter{Analyse et conception}
\label{ch:conception}

Ce chapitre est le c\oe ur du rapport. Il fixe les besoins, le modèle de menace et
l'invariant, puis présente l'architecture --- ce que l'outil lit, voit, refuse, demande et
enregistre. Il expose ensuite les principes qui gouvernent chaque décision, avec, pour
chacun, la commande qui casse l'approche naïve et le test qui fige la règle. Il décrit
enfin la conception du moteur, de la politique et du décideur unique, la répartition entre
Linux et Windows, et résume les choix d'architecture face à leurs alternatives.

% 3.1 — Besoins
\section{Analyse des besoins}
\label{sec:besoins}

\subsection{Acteurs}
Trois acteurs interagissent avec \demo{} : l'\textbf{agent}, qui propose des commandes et
des modifications de fichiers ; la \textbf{personne} (le développeur), qui installe l'outil,
répond aux questions qu'il pose et restaure ce qu'il a capturé ; et les \textbf{systèmes
tiers} touchés par les commandes (système de fichiers local, bases de données, dépôts git
distants, fournisseurs cloud).

\subsection{Besoins fonctionnels}
Le tableau~\ref{tab:bf} rassemble ce que l'outil doit faire pour ces acteurs.

\begin{table}[htbp]
\centering
\caption{Besoins fonctionnels}
\label{tab:bf}
\small
\begin{tabularx}{\linewidth}{@{}L{1cm}Y@{}}
\toprule
\textbf{Id} & \textbf{Besoin} \\
\midrule
BF1 & Intercepter, avant exécution, chaque commande shell (POSIX et PowerShell) et chaque
      modification de fichier proposée par l'agent, par toutes les portes disponibles. \\
BF2 & Prédire l'effet de la commande : ce qu'elle détruit, crée, déplace, et ce qu'elle
      atteint à distance. \\
BF3 & Rendre une décision parmi six : ALLOW, REVERSIBLE, DRY\_RUN, SUPERVISED,
      CONTEXT\_MISMATCH, ESCALATE ; pour une escalade, indiquer si une personne peut
      l'approuver sur le champ. \\
BF4 & Capturer, avant l'exécution, une copie exacte de ce qui sera détruit, en un seul
      groupe par commande, restaurable par \cmd{demo\_cli undo <id>}. \\
BF5 & Prévisualiser une requête SQL destructrice (nombre de lignes touchées). \\
BF6 & Placer sous un filet d'exécution ce qui ne peut pas être lu (Linux et Windows). \\
BF7 & Refuser toute modification des fichiers de \demo{} lui-même par l'agent. \\
BF8 & Enregistrer chaque décision dans un registre infalsifiable, vérifiable par
      \cmd{demo\_cli verify}. \\
BF9 & Diagnostiquer l'installation (\cmd{demo\_cli doctor}) et dire, par agent et par
      système, ce qui est protégé et ce qui ne l'est pas. \\
\bottomrule
\end{tabularx}
\end{table}

\subsection{Besoins non fonctionnels}
Le tableau~\ref{tab:bnf} fixe les qualités attendues, chacune avec sa justification.

\begin{table}[htbp]
\centering
\caption{Besoins non fonctionnels}
\label{tab:bnf}
\small
\begin{tabularx}{\linewidth}{@{}L{1.2cm}L{3.4cm}Y@{}}
\toprule
\textbf{Id} & \textbf{Qualité} & \textbf{Exigence et justification} \\
\midrule
BNF1 & Honnêteté & Une décision ne revendique jamais une protection non prise ; la raison
  dit ce qui a été capturé et ce qui ne l'a pas été. \\
BNF2 & Sûreté par défaut & Une \emph{décision} incertaine va vers le plus strict ; une
  ligne non lue n'est jamais ALLOW. \\
BNF3 & Ne pas bloquer l'agent sur une panne & Une erreur interne de \demo{} ne doit pas
  paralyser l'agent : elle laisse passer, bruyamment ; mais une panne du \emph{moteur} ne
  rend jamais une décision plus permissive (B0). \\
BNF4 & Latence & Le moteur dispose de 10\,ms de temps CPU (50\,ms pour une seconde
  tentative) et d'un plafond de 250\,ms en temps réel ; les familles d'outils ne se chargent
  qu'à la demande. Un outil lent est un outil désinstallé. \\
BNF5 & Portabilité & Linux et Windows (PowerShell, Git Bash), avec les mêmes garanties ;
  chaque partie est vérifiée sur les deux systèmes. \\
BNF6 & Non-régression & Aucune décision ne devient plus permissive sans arbitrage
  explicite (portail de \nbcorpus{} commandes). \\
BNF7 & Confidentialité & Les reçus masquent les secrets (mots de passe d'URL, clés) ;
  aucune télémétrie. \\
\bottomrule
\end{tabularx}
\end{table}

% 3.2 — Modèle de menace et invariant
\section{Modèle de menace et invariant}
\label{sec:menace}

\subsection{Contre quoi \texorpdfstring{\demo}{demo\_cli} protège}
Le modèle de menace est écrit dans le code (\fichier{protection.py}) et affiché par
\cmd{demo\_cli doctor} :

\begin{encadre}[Modèle de menace (\fichier{protection.py})]
\textit{« demo\_cli protects against an agent's MISTAKES, not a determined adversary:
commands it understands are judged before they run; for commands it cannot read, the
runtime net (syscall guard on Linux, WinFsp mount on Windows) is the floor; a remote change
that can lose data or cannot be recalled is stopped for review; one that only adds [\ldots]
is allowed and listed. »}
\end{encadre}

L'adversaire visé est donc un agent \emph{coopératif qui se trompe} : il veut nettoyer un
dossier et en supprime un autre, il croit qu'un répertoire existe, il réinitialise une base
qu'il croit de test. Ce choix a deux conséquences assumées.
\begin{itemize}
  \item Ce qui ne relève pas d'une erreur plausible mais d'une évasion délibérée (une charge
    chiffrée, une course entre la vérification et l'usage d'un appel système) n'est pas
    l'objectif premier ; le filet d'exécution en réduit cependant la portée.
  \item En revanche, l'agent ne doit pas pouvoir \emph{désarmer} l'outil par accident ou par
    excès de zèle : modifier la configuration de \demo{}, désinstaller son \textit{hook},
    effacer son registre. Ces actions sont refusées à toutes les portes
    (section~\ref{sec:plan-controle}).
\end{itemize}

\subsection{L'invariant}
Toute la logique de décision découle d'une phrase, héritée de la bêta et gravée dans chaque
reçu (champ \cmd{invariant}) :

\begin{encadre}[Invariant]
\texttt{mutating\_actions\_must\_be\_recoverable\_and\_match\_the\_declared\_context}\\[2pt]
Une action qui modifie doit être récupérable \emph{et} conforme au contexte déclaré ;
sinon, elle est escaladée.
\end{encadre}

\subsection{Fermé pour les décisions, ouvert pour les pannes}
Un outil de sécurité branché sur un agent fait face à une tension : s'il échoue « fermé »,
le moindre défaut interne bloque tout le travail de l'agent et l'outil est désinstallé ;
s'il échoue « ouvert », un défaut devient une brèche. \demo{} tranche ainsi :
\begin{itemize}
  \item une \textbf{décision} est toujours fermée : sans point de récupération, une
    modification est escaladée dans \emph{tout} environnement, développement compris ;
  \item une \textbf{erreur interne} de l'adaptateur (entrée illisible, exception) laisse
    passer la commande, mais \emph{bruyamment} : un avertissement visible, jamais un
    silence ;
  \item une \textbf{défaillance du moteur} (dépassement de son budget, exception) ne rend
    jamais une décision plus permissive : l'ancien chemin peut rester plus strict, mais son
    ALLOW devient \cmd{engine\_unavailable}, puis filet, point de contrôle ou escalade
    (partie B0).
\end{itemize}

% 3.3 — Architecture d'ensemble
\section{Architecture d'ensemble}
\label{sec:architecture}

\subsection{Le flux d'une commande}
La figure~\ref{fig:flux} suit une commande de l'agent jusqu'au reçu. Cette section présente
l'architecture par ses idées ; l'annexe~\ref{ann:architecture} la montre en entier, module par
module.

\begin{figure}[htbp]
\centering
% Figure 1 — Le flux d'ensemble : de l'agent au reçu.
\begin{tikzpicture}[node distance=5mm and 6mm]
  % --- Rangée du haut : agent -> portes -> moteur
  \node[douce,text width=26mm] (agent) at (0,0) {Agent de code IA\sous{Claude Code, Codex}};
  \node[boite,text width=30mm] (edit) at (4.6,0) {Porte fichier\sous{Edit, Write}};
  \node[boite,text width=30mm,above=2mm of edit] (hook) {\textit{Hook} PreToolUse\sous{Bash, PowerShell}};
  \node[boite,text width=30mm,below=2mm of edit] (shell) {Garde du shell\sous{piège bash, \texttt{check}}};
  \node[zone,fit=(hook)(edit)(shell),label={[etiquette]above:portes}] (portes) {};
  \draw[fleche] (agent.east) -- (portes.west |- agent.east);

  \node[accent,text width=36mm] (sim) at (9.9,0) {Simulation\sous{FS virtuel (\texttt{simulate.py})}};
  \node[accent,text width=36mm,above=3mm of sim] (parse) {Analyse syntaxique\sous{POSIX \textbar{} PowerShell (C3)}};
  \node[accent,text width=36mm,below=3mm of sim] (eff) {Effets $\rightarrow$ impact\sous{détruit, crée, distant, inconnu}};
  \node[zone,draw=wepwnfonce,fit=(parse)(sim)(eff),label={[etiquette,text=wepwnfonce]above:moteur}] (moteur) {};
  \draw[fleche] (parse) -- (sim);
  \draw[fleche] (sim) -- (eff);
  \draw[fleche] (portes.east |- sim.west) -- (moteur.west |- sim.west);

  % --- Politique et décideur
  \node[boite,text width=36mm] (pol) at (9.9,-3.4) {\texttt{verdict} + \texttt{policy}\sous{effet $\times$ surface}};
  \node[boite,draw=encre,line width=1.1pt,text width=30mm] (dec) at (4.6,-3.4) {\textbf{\texttt{decide.py}}\sous{le décideur unique}};
  \draw[fleche] (moteur.south) -- (pol.north);
  \draw[fleche] (pol.west) -- (dec.east);

  % --- Décisions (colonnes fixes)
  \foreach \x/\n/\c/\t in {0/o1/dALLOW/ALLOW,
                          2.6/o2/dREV/{REVERSIBLE\\[-2pt]{\tiny ou DRY\_RUN}},
                          5.2/o3/dSUP/SUPERVISED,
                          7.8/o4/dCTX/{CONTEXT\\[-2pt]MISMATCH},
                          10.4/o5/dASK/{ESCALATE\\[-2pt]{\tiny ask ou deny}}}
    \node[decision,fill=\c,text width=21mm] (\n) at (\x,-5.4) {\t};
  \draw[fleche] (dec.south) -- ++(0,-0.45) coordinate (b);
  \draw[encre!70,line width=0.7pt] (o1.north |- b) -- (o5.north |- b);
  \foreach \n in {o1,o2,o3,o4,o5} \draw[fleche] (\n.north |- b) -- (\n.north);

  % --- Conséquences
  \node[douce,text width=21mm] (rec) at (2.6,-7.0) {Point de\\récupération\sous{\texttt{undo <id>}}};
  \node[douce,text width=21mm] (net) at (5.2,-7.0) {Filet\\d'exécution\sous{ptrace / WinFsp}};
  \node[douce,text width=21mm] (per) at (9.1,-7.0) {La personne\\tranche\sous{ou refus net}};
  \draw[fleche,densely dashed] (o2) -- (rec);
  \draw[fleche,densely dashed] (o3) -- (net);
  \draw[fleche,densely dashed] (o4.south) -- (o4.south |- per.north);
  \draw[fleche,densely dashed] (o5.south) -- (o5.south |- per.north);

  % --- Reçu
  \node[boite,draw=wepwnfonce,line width=0.9pt,minimum width=128mm] (recu) at (5.2,-8.7)
       {Reçu dans le registre chaîné (SHA-256) --- \emph{chaque} décision en laisse un};
  \foreach \n in {o1,rec,net} \draw[fleche] (\n.south) -- (\n.south |- recu.north);
  \draw[fleche] (per.south) -- (per.south |- recu.north);
\end{tikzpicture}

\caption{Le flux d'ensemble : de la commande de l'agent au reçu}
\label{fig:flux}
\end{figure}

\begin{enumerate}
  \item \textbf{La porte} reçoit l'action : le \textit{hook} PreToolUse de Claude Code (outils
    \cmd{Bash} et \cmd{PowerShell}), la porte fichier (\cmd{Edit}, \cmd{Write}), le
    \textit{hook} de Codex ou de Cursor, le piège du shell pour les commandes tapées dans un
    terminal, ou \cmd{demo\_cli check}. Chaque porte ne fait que traduire : elle passe la
    commande à l'orchestrateur et traduit la décision dans ce que son agent comprend.
  \item \textbf{L'orchestrateur} (\fichier{guard.py}, \cmd{Guard.evaluate}) enchaîne des
    étapes nommées : refus d'une configuration cassée ; prédiction du moteur ; refus de ce
    qui touche le plan de contrôle ; choix du chemin (moteur, ligne non lue, ancien chemin) ;
    contexte ; décision ; reçu.
  \item \textbf{Le moteur} analyse la ligne, la simule sur une vue virtuelle du système de
    fichiers et produit un \emph{impact} : ce qui est détruit, créé, atteint à distance, non
    compris, non couvert.
  \item \textbf{La politique} (\fichier{policy.py}) traduit chaque effet distant en
    décision, chaque étiquette en surface, et chaque surface en canal « demander » ou
    « refuser ».
  \item \textbf{Le décideur} (\fichier{decide.py}), fonction pure sans entrée-sortie, rend la
    décision. Il ne capture rien lui-même : il reçoit seulement le fait qu'une capture a
    \emph{réussi}. C'est ce qui l'empêche de mentir.
  \item \textbf{Les conséquences} : une capture (REVERSIBLE), un filet (SUPERVISED), une
    question à la personne ou un refus (ESCALATE) --- et, dans tous les cas, un reçu.
\end{enumerate}

\subsection{Ce que l'outil voit}
\label{sec:vision}
La figure~\ref{fig:vision} répond à la question qu'un jury pose en premier : \emph{que
voit réellement l'outil, et où s'arrête-t-il ?} Elle croise six couches --- lire,
récupérer, surveiller, refuser, demander, tracer --- avec trois zones où une commande peut
agir.

\begin{figure}[htbp]
\centering
% Figure 2 — Ce que l'outil voit : chaque couche, et ce qu'elle couvre (ou pas),
% selon l'endroit où la commande agit.
\begin{tikzpicture}[
    cell/.style={draw=white,line width=1.2pt,minimum width=38mm,minimum height=12.5mm,
                 text width=35mm,align=center,font=\sffamily\scriptsize,inner sep=2pt},
    oui/.style={cell,fill=wepwnclair},
    part/.style={cell,fill=dASK!18},
    non/.style={cell,fill=grisclair,text=gris,
                postaction={pattern=north east lines,pattern color=trait}},
    ligne/.style={minimum width=34mm,minimum height=12.5mm,text width=32mm,
                  align=left,font=\sffamily\small,inner sep=2pt},
    tete/.style={minimum width=38mm,font=\sffamily\small\bfseries,align=center,text width=36mm},
  ]
  \def\W{3.85}\def\H{1.3}
  % En-têtes de colonnes
  \node[tete] at (1*\W,0.2) {Dans le projet};
  \node[tete] at (2*\W,0.2) {Sur la machine,\\hors du projet};
  \node[tete] at (3*\W,0.2) {Surfaces distantes\\{\scriptsize\mdseries base, cloud, git, SaaS}};

  % Lignes : (y, titre, sous-titre)
  \foreach \y/\t/\s in {1/Lire avant/le moteur (statique),
                        2/Récupérer/{instantané, point de contrôle},
                        3/Surveiller/{filet : ptrace ou WinFsp},
                        4/Refuser/plan de contrôle,
                        5/Demander/la personne (ask),
                        6/Tracer/les reçus}
    \node[ligne] at (0.15,-\y*\H) {\textbf{\t}\\{\scriptsize\color{gris}\s}};

  % Lire
  \node[oui]  at (1*\W,-1*\H) {prédit ce qui est détruit,\\créé, déplacé};
  \node[oui]  at (2*\W,-1*\H) {prédit la destruction\\$\rightarrow$ ESCALATE};
  \node[oui]  at (3*\W,-1*\H) {reconnaît la surface\\et la sorte d'effet};
  % Récupérer
  \node[oui]  at (1*\W,-2*\H) {copie exacte avant,\\tout ou rien};
  \node[non]  at (2*\W,-2*\H) {aucune copie\\(V17 : jamais prétendue)};
  \node[non]  at (3*\W,-2*\H) {aucune copie locale ;\\indice chez le fournisseur};
  % Surveiller
  \node[oui]  at (1*\W,-3*\H) {sauvegarde chaque fichier\\au moment où il est détruit};
  \node[part] at (2*\W,-3*\H) {enregistré, non sauvegardé\\(V18, filet Linux)};
  \node[non]  at (3*\W,-3*\H) {invisible au filet\\local};
  % Refuser
  \node[oui]  at (1*\W,-4*\H) {\texttt{.demo\_cli/}, \texttt{.demo\_cli.toml},\\\textit{hooks} du projet};
  \node[part] at (2*\W,-4*\H) {chemins listés seulement :\\\texttt{\textasciitilde/.bashrc}, \textit{hooks}, tâche};
  \node[non]  at (3*\W,-4*\H) {sans objet};
  % Demander
  \node[part] at (1*\W,-5*\H) {seulement si rien\\n'a pu être capturé};
  \node[part] at (2*\W,-5*\H) {destruction hors projet\\: ask};
  \node[oui]  at (3*\W,-5*\H) {ask ou deny\\selon la surface};
  % Tracer
  \node[oui]  at (1*\W,-6*\H) {chaque décision};
  \node[oui]  at (2*\W,-6*\H) {chaque décision};
  \node[oui]  at (3*\W,-6*\H) {chaque décision\\(sans secrets)};

  % Légende
  \node[oui,minimum width=6mm,text width=4mm,minimum height=4mm] (l1) at (0.4,-8.75) {};
  \node[right=1mm of l1,font=\sffamily\scriptsize] (t1) {couvert};
  \node[part,minimum width=6mm,text width=4mm,minimum height=4mm,right=5mm of t1] (l2) {};
  \node[right=1mm of l2,font=\sffamily\scriptsize] (t2) {partiel};
  \node[non,minimum width=6mm,text width=4mm,minimum height=4mm,right=5mm of t2] (l3) {};
  \node[right=1mm of l3,font=\sffamily\scriptsize] {non couvert --- dit tel quel dans la raison de la décision};
\end{tikzpicture}

\caption{Ce que l'outil voit : chaque couche et ce qu'elle couvre, selon l'endroit où la commande agit}
\label{fig:vision}
\end{figure}

Trois lectures de cette figure méritent d'être soulignées.
\begin{itemize}
  \item \textbf{Seule la lecture voit partout} --- mais seulement ce qu'elle sait lire. Tout
    le reste du dispositif existe pour ce qu'elle ne sait pas lire.
  \item \textbf{La récupération et le filet s'arrêtent au projet.} Hors du projet, aucune copie
    n'est possible, et la décision le dit ; une destruction connue hors projet est escaladée, et
    celle que fait du code non lu est seulement enregistrée par le filet Linux (V18).
  \item \textbf{Le distant ne se récupère pas, il se décide.} Pour une base ou un dépôt
    distant, \demo{} ne peut que reconnaître la surface, demander ou refuser, et indiquer où
    le fournisseur garde son propre historique.
\end{itemize}

% 3.4 — Les principes qui gouvernent les décisions
\section{Les principes qui gouvernent les décisions}
\label{sec:principes}

Chaque principe ci-dessous est présenté de la même façon : l'énoncé, la commande qui casse
l'approche naïve, la règle retenue et le test qui la fige. Les sorties de terminal sont
réelles (commit \commitref{}, Linux, sans filet d'exécution sauf mention contraire) ; seules
les lignes sans intérêt sont coupées et les chemins abrégés en \cmd{\textasciitilde/projet}.

% ------------------------------------------------------------------
\subsection{Analyser et simuler plutôt que reconnaître des motifs}
\label{sec:pourquoi-moteur}

\paragraph{Le constat.} La bêta et la version 1.7.0 lisaient les commandes avec un
découpeur de texte et des expressions régulières. L'audit et la revue menés au début de la
période documentée ont trouvé, encore et encore, des défauts de la même forme : \emph{deux
analyseurs répondant différemment à la même question} --- le découpeur de lignes, le
normaliseur de mot de commande, le masqueur de citations, le lecteur de \textit{heredocs},
le découpeur d'opérandes. Deux de ces défauts avaient été introduits par des corrections
des autres (\fichier{tests/test\_command\_model.py}).

\paragraph{Ce qui casse.} Le shell n'est pas un texte mais une grammaire. Un commentaire
neutralise ce qui le suit (\cmd{ls \# ; rm -rf x} ne supprime rien) ; une substitution
\cmd{\$(\ldots)} exécute une commande cachée dans un argument (\cmd{git commit -m "\$(rm
y)"}) ; une enveloppe déplace le vrai mot de commande (\cmd{sudo -u x rm},
\cmd{bash -o pipefail -c "rm app.db"}, \cmd{env -S 'rm -rf x'}) ; un \cmd{case} ou une
fonction contiennent des commandes que le découpage par \cmd{;} ne voit pas ; une citation
ANSI-C (\cmd{\$'\textbackslash x72m'}) écrit \cmd{rm} sans les lettres \cmd{r} et \cmd{m}.

\paragraph{La règle retenue.} Une seule analyse de la ligne (\fichier{model.py}), lue par
tout le reste. Avant de remplacer le découpeur, elle a tourné en ombre : sur toutes les
commandes de la suite de tests (801, dans les deux dialectes), elle a rendu la même décision
que le découpeur ; sur un jeu étiqueté à la main de ce que bash exécute réellement, elle
avait raison dans 22 cas où le découpeur avait tort, et tort dans aucun cas où il avait
raison. Ces deux chiffres justifient le remplacement plutôt que la cohabitation.

\paragraph{Et la simulation ?} Une structure correcte ne suffit pas
(section~\ref{sec:statique}) : \cmd{mv a.txt b.txt} n'est destructeur que si \cmd{b.txt}
existe. Le moteur rejoue donc la ligne sur une vue virtuelle du système de fichiers et suit
l'origine des contenus : dans \cmd{mv a b \&\& rm b}, ce qui est perdu est le contenu de
\cmd{a}. C'est exactement le raisonnement qui manquait à l'agent dans l'incident Gemini CLI.

% ------------------------------------------------------------------
\subsection{L'inconnu n'est pas sûr : « non couvert » et « non compris »}
\label{sec:inconnu}

\paragraph{Deux ignorances différentes.} Le moteur distingue :
\begin{itemize}
  \item \textbf{non couvert} (\textit{uncovered}) : la ligne est lue, mais un outil n'a pas
    de règle. On sait \emph{ce qui} est appelé, avec quels arguments ;
  \item \textbf{non compris} (\textit{not understood}) : une partie de la ligne ne peut pas
    être lue avant l'exécution --- une variable affectée ailleurs, une substitution dont la
    sortie n'est pas prévisible, un mot de commande construit.
\end{itemize}

\paragraph{La règle retenue (parties M2.3a et M2.3b, arbitrages M1--M5 et Rb1--Rb4).}
\begin{enumerate}
  \item Une ligne \textbf{non comprise} va au filet d'exécution s'il est disponible, sinon à
    un point de contrôle s'ils sont activés, sinon elle est \textbf{escaladée}. Elle
    n'obtient jamais ALLOW.
  \item Un \textbf{exécuteur de code} (\cmd{python}, \cmd{node}, \cmd{make}, \cmd{npm},
    \cmd{./script.sh}, \cmd{source}\ldots) lance du code que \demo{} ne lit pas : il va au
    filet, sinon au point de contrôle, sinon il est autorisé \emph{avec une note qui dit
    qu'aucune capture n'a été faite} (arbitrage M1). Il garde le filet même si on lui passe
    un fichier hors du projet (V21).
  \item Deux cas d'exécuteur sont plus risqués et \textbf{escaladent sans filet} (supervisés
    avec, arbitrage Rb2) : du \textbf{code en ligne qui lance d'autres programmes}
    (\cmd{python -c "\ldots os.system(\ldots)"}, \cmd{node -e "\ldots child\_process\ldots"})
    et les \textbf{installations de paquets} (\cmd{pip install}, \cmd{npm install}\ldots), qui
    exécutent les scripts d'installation de chaque paquet.
  \item Une ligne qui atteint quelque chose au-delà des fichiers locaux sans être lue
    (cloud, conteneurs, services, une option inconnue sur un outil distant, \cmd{kubectl
    exec}, \cmd{heroku run}\ldots) est escaladée : ni le filet ni le point de contrôle ne la
    voient. Le code récupéré du réseau puis exécuté (\cmd{curl \ldots{} | bash}) garde son
    refus.
  \item Un mot sans règle qui n'est pas un exécuteur de code n'a pas d'effet (arbitrage Rb1).
\end{enumerate}

\begin{terminal}[ligne non comprise, pas de filet]
$ demo_cli check --mode enforce 'rm $INCONNU'
demo_cli 1.7.0  [x] BLOCKED  ·  ESCALATE
  matched rule        not_understood
Decision
  No recovery path for a not_understood on 'unknown'; cannot auto-recover.
  Human input required.
$ demo_cli check --mode enforce 'python3 outil.py'
demo_cli 1.7.0  [+] SAFE  ·  ALLOW
Decision
  Allowed: it runs code demo_cli does not read (python3); no runtime net was
  available, so nothing that code changes is captured or watched.
$ demo_cli check --mode enforce 'pip install requests'
demo_cli 1.7.0  [x] BLOCKED  ·  ESCALATE
Decision
  No recovery path for a unread_code on 'unknown'; cannot auto-recover. Human input
  required. It installs packages, which run their own install scripts (pip), and no
  runtime net was available to capture what that changes.
\end{terminal}

\begin{pourquoi}[Pourquoi l'exécuteur sans filet est-il autorisé ?]
Escalader chaque \cmd{python script.py}, \cmd{make} ou \cmd{npm test} rendrait l'outil
inutilisable : c'est l'essentiel du travail d'un agent. L'arbitrage M1 accepte ce cas \emph{à
visage découvert} : la raison de la décision dit qu'aucune capture n'a été faite. Les deux cas
où l'exécuteur lance à coup sûr du code étranger --- un programme démarré depuis du code en
ligne, une installation de paquets --- escaladent en l'absence de filet (Rb2). Là où le filet
existe (Linux avec Claude Code, Windows avec le montage), ces lignes sont SUPERVISED. Un
outil inconnu qui n'est pas un exécuteur de code (« non couvert ») est, lui, traité comme
sans effet (Rb1) ; la table de règles s'étend famille par famille pour réduire ce cas, et un
test (\fichier{test\_m22f.py}) vérifie qu'aucune ligne du portail sur laquelle une ancienne
règle se déclenchait n'est restée hors du moteur sans être marquée non comprise.
\end{pourquoi}

\paragraph{Lire plutôt qu'abandonner.} « Non compris » doit rester rare, sinon tout escalade.
Le moteur lit donc ce qui peut l'être sans rien exécuter : la sortie de \cmd{\$(cat
liste.txt)}, \cmd{echo}, \cmd{printf}, \cmd{pwd}, \cmd{basename}, \cmd{dirname}, sous des
conditions strictes (fichier régulier du projet, UTF-8, au plus 64~Kio, non modifié plus tôt
sur la ligne ; partie I1b). La partie M2.3b a ajouté ce que l'ancien chemin savait capturer :
une valeur connue sur la ligne lue comme l'option qu'elle est (\cmd{X=-rf; rm \$X src}), un
script littéral envoyé à un shell par un tube (\cmd{echo 'rm -rf build' | sh}),
\cmd{sqlite3 \$DB} et \cmd{sqlite3 app.db "\$(cat q.sql)"}, les messages de commit, et
\cmd{npx}, \cmd{bunx} ou \cmd{pnpm dlx} qui lancent un outil connu (\cmd{npx prisma migrate
reset} est lu comme \cmd{prisma}) :

\begin{terminal}[sortie de substitution lue par le moteur]
$ demo_cli check --mode enforce 'rm $(cat list.txt)'
demo_cli 1.7.0  [+] SAFE  ·  REVERSIBLE
  target              ~/projet/a.txt
Recovery
  snapshot            a.txt.20261009-175242.6d6ca4f4.bak
Decision
  Recovery point captured before the action; reversible.
\end{terminal}

% ------------------------------------------------------------------
\subsection{Jamais un REVERSIBLE non mérité}
\label{sec:jamais-reversible}

\paragraph{L'énoncé.} REVERSIBLE est une promesse : « \cmd{demo\_cli undo} remettra tout en
place ». Une capture partielle, ou une copie qui ne contient pas ce que la commande détruit,
rendrait cette promesse fausse, ce qui est pire qu'une escalade : la personne cesse de se
méfier.

\paragraph{Les règles qui en découlent.}
\begin{itemize}
  \item \textbf{Tout ou rien} : les captures d'une commande forment un groupe ; si l'une échoue,
    les autres sont invalidées et la commande est escaladée (partie S2,
    \cmd{recovery.snapshot\_set}).
  \item \textbf{Pas de point de contrôle pour ce qu'il ne contient pas} : une copie du projet
    ne restaure ni un fichier hors du projet (V17), ni une base distante, ni un historique
    git poussé de force.
  \item \textbf{Hors projet, rien n'est promis} : une destruction connue hors du projet est
    escaladée, même à côté d'une partie non comprise.
  \item \textbf{Le décideur ne capture pas} : \cmd{decide.py} ne reçoit qu'un booléen « une
    capture de la cible réelle a réussi », fourni par l'orchestrateur ; il n'existe aucune
    cible par défaut sur laquelle se rabattre.
\end{itemize}

% ------------------------------------------------------------------
\subsection{Pendant la migration, la réponse la plus stricte gagne}
\label{sec:plus-strict}

\paragraph{Le problème.} Le moteur a remplacé l'ancien chemin \emph{progressivement}, famille
par famille. Pendant ce temps, deux lecteurs coexistent : le moteur et les anciennes règles.
Que faire quand ils divergent ?

\paragraph{La règle retenue (parties B0 et M2.3a).} L'ancien chemin peut garder une
\emph{capture} ou une \emph{escalade}, jamais son ALLOW. Exemple réel, au moment de M2.3a : le
moteur marquait \cmd{X=-rf; rm \$X build} comme non compris (une valeur qui commence par un
tiret) ; l'ancienne règle, elle, savait capturer \cmd{build}. La réponse la plus stricte ---
capturer --- l'emportait. À l'inverse, quand l'ancienne règle aurait autorisé une ligne que le
moteur ne lit pas, c'est la règle M3 qui s'applique (filet, point de contrôle, escalade) : la
vulnérabilité V19 était précisément un \cmd{psql -c "\$(cat q.sql)"} autorisé alors que
\cmd{q.sql} contenait \cmd{DROP TABLE users;}. Il est désormais escaladé.

\paragraph{La fin de la migration (M2.3b).} La partie M2.3b a fait lire au moteur tout ce que
l'ancien chemin capturait encore ; depuis, \textbf{l'ancien chemin ne décide plus aucune ligne
POSIX}. La même ligne est maintenant capturée par le moteur lui-même :

\begin{terminal}[la ligne est lue par le moteur (M2.3b)]
$ demo_cli check --mode enforce 'X=-rf; rm $X build'
demo_cli 1.7.0  [+] SAFE  ·  REVERSIBLE
  target              ~/projet/build
Recovery
  snapshot            build.20261009-175242.30656068.snapdir
Decision
  Recovery point captured before the action; reversible.
\end{terminal}

La règle « le plus strict gagne » ne vaut plus que là où l'ancien chemin sert encore : le texte
\cmd{cmd.exe}, une ligne POSIX qui nomme une cmdlet PowerShell (M5), et le cas où le moteur
échoue (B0). Une construction PowerShell que C3 ne lit pas
n'y retourne pas : elle est non comprise, et suit le même chemin qu'une ligne POSIX non lue
--- filet, sinon point de contrôle, sinon escalade.

% ------------------------------------------------------------------
\subsection{Les fichiers de \texorpdfstring{\demo}{demo\_cli} : le plan de contrôle}
\label{sec:plan-controle}

\paragraph{L'énoncé.} Un agent qui peut modifier la configuration de son garde-fou n'a pas de
garde-fou. Le \emph{plan de contrôle} regroupe : la configuration \fichier{.demo\_cli.toml},
l'espace de travail \fichier{.demo\_cli/} (registre, réserve de récupération), les
configurations des \textit{hooks} dans le projet et dans le répertoire personnel, le script du
garde du shell, \fichier{\textasciitilde/.bashrc} quand il l'active, le lanceur
\cmd{demo\_cli}, et sous Windows les fichiers de la tâche d'ouverture de session.

\paragraph{La règle.} Toute modification de ces chemins --- y compris une \emph{création} ou un
ajout en fin de fichier --- est refusée (ESCALATE, jamais approuvable sur place), à toutes les
portes : avant l'exécution par le moteur et par les anciennes règles, dans la porte fichier,
et pendant l'exécution par le garde d'appels système, qui annule l'appel.

\begin{terminal}[un agent tente de passer l'outil en mode observation]
$ demo_cli check --mode enforce 'echo mode = "shadow" > .demo_cli.toml'
demo_cli 1.7.0  [x] BLOCKED  ·  ESCALATE
  matched rule        control_plane
Decision
  This changes demo_cli's own files, or a directory holding them (.demo_cli.toml):
  its config, hook settings, ledger or recovery store. An agent may not [...]
\end{terminal}

% ------------------------------------------------------------------
\subsection{Demander ou refuser}
\label{sec:ask-deny}

\paragraph{Le problème.} Une escalade arrête l'agent. Faut-il alors \emph{demander} à la
personne, ou \emph{refuser} ? Demander tout use l'attention (section~\ref{sec:permissions}) ;
refuser tout rend l'outil inutilisable pour des opérations ordinaires (supprimer un volume
Docker de test, pousser une branche).

\paragraph{La règle retenue (partie ASK).} L'escalade reste une escalade --- aucune décision du
portail ne change --- mais chaque surface porte un \emph{canal} dans
\fichier{policy.ESCALATION\_CHANNEL} :
\begin{itemize}
  \item \textbf{ask} quand le changement est réel mais ordinaire, et que la commande exacte suffit
    à la personne pour savoir à quoi elle consent (base distante, historique git, bucket,
    conteneur, déploiement, service\ldots) ;
  \item \textbf{deny} quand tout part sans retour réaliste (infrastructure détruite, vidage d'un
    entrepôt de données, formatage de disque, paiement, compte utilisateur supprimé, données
    de récupération du système) ; seuls un jeton d'approbation structurel, que l'agent ne
    peut pas forger, ou une exécution hors de l'agent le permettent.
\end{itemize}
\textbf{Une surface sans ligne dans la table est refusée} : une surface ajoutée plus tard ne
devient jamais « demandable » par accident. Enfin, \emph{demander} n'a de sens que si
quelqu'un voit la question : le \textit{hook} Claude Code ne pose la question que dans les
modes où l'interface l'affiche (\cmd{default}, \cmd{acceptEdits}, \cmd{plan}, \cmd{auto}) et
refuse dans les autres --- une question que personne ne voit ne doit jamais devenir « c'est
passé ». Codex, qui ne sait pas demander, refuse.

% ------------------------------------------------------------------
\subsection{Les faits sur les outils sont vérifiés}
\label{sec:faits-verifies}

Une règle fausse est une décision fausse. Chaque grammaire d'options est tirée de la source
de l'outil lui-même --- son \cmd{--help}, sa page de manuel, son code source, ses données
publiées (l'index des commandes d'\cmd{aws}, l'arbre de commandes de \cmd{gcloud}, la table
de commandes de \cmd{redis}) --- et vérifiée sur la version courante. Les comportements
douteux sont testés sur des fichiers jetables : que \cmd{unzip -Pxo} lit \cmd{xo} comme mot de
passe et non comme l'option d'écrasement (I1c, vérifié sur UnZip 6.0), que \cmd{helm 4}
change la sémantique de \cmd{rollback --dry-run}, que \cmd{tofu} accepte \cmd{-chdir DIR}
sans \cmd{=} quand \cmd{terraform} le refuse. Un test confronte chaque grammaire au
\cmd{--help} de l'outil installé. Une règle dont la grammaire vient seulement de la
documentation porte \cmd{verified=False} : c'est le cas de \nbnonverifiees{} règles sur
\nbregles{}.

% 3.5 — Conception du moteur
\section{Conception du moteur}
\label{sec:moteur}

\subsection{Vue d'ensemble}
Le moteur (\fichier{src/demo\_cli/engine/}) transforme une ligne en un \emph{impact}, sans
rien exécuter ni rien écrire. La figure~\ref{fig:moteur} le montre sur un exemple réel.

\begin{figure}[htbp]
\centering
\makebox[\linewidth][c]{% Figure 3 — L'intérieur du moteur, sur un exemple réel (sortie du moteur au commit de référence).
\begin{tikzpicture}[
    etape/.style={accent,text width=31mm,minimum height=11mm},
    ex/.style={douce,text width=82mm,align=left,font=\ttfamily\scriptsize,minimum height=11mm,inner sep=4pt},
    lien/.style={densely dotted,draw=gris,line width=0.6pt},
  ]
  \def\D{1.62}
  \node[douce,text width=31mm,minimum height=9mm] (l) at (0,0) {Ligne de l'agent};
  \node[ex,minimum height=9mm] (el) at (6.9,0) {\color{encre}\normalsize\ttfamily X=build; rm -rf \$X \&\& mv a.txt b.txt};

  \node[etape] (p) at (0,-1*\D) {Analyse syntaxique\\{\scriptsize\texttt{model.py}}};
  \node[ex] (ep) at (6.9,-1*\D) {3 commandes : [X=build] ; [rm -rf \$X] \&\& [mv a.txt b.txt]\\
     \sffamily mots avec leurs citations, \texttt{\$X} marqué « s'étend »};

  \node[etape] (i) at (0,-2*\D) {Invocations\sous{grammaire de l'outil}};
  \node[ex] (ei) at (6.9,-2*\D) {rm : options \{force, recursive\}, opérandes [\$X]\\
     mv : opérandes [a.txt, b.txt]};

  \node[etape] (w) at (0,-3*\D) {Expansion~des~mots\\{\scriptsize\texttt{words.py}}};
  \node[ex] (ew) at (6.9,-3*\D) {\$X $\rightarrow$ build \sffamily(affectation plus tôt sur la ligne)\\
     \sffamily inconnu $\Rightarrow$ « non compris », jamais deviné};

  \node[etape] (s) at (0,-4*\D) {Simulation\sous{FS virtuel, \texttt{simulate.py}}};
  \node[ex] (es) at (6.9,-4*\D) {rm -r : DELETE build/ \sffamily(existe)\\
     \ttfamily mv : b.txt existe $\rightarrow$ OVERWRITE b.txt \sffamily(origine suivie)};

  \node[etape,draw=encre,line width=0.9pt] (m) at (0,-5*\D) {Impact\\{\scriptsize\texttt{effects.Impact}}};
  \node[ex,fill=white,draw=wepwnfonce] (em) at (6.9,-5*\D) {destroys = \{build: DELETE, b.txt: OVERWRITE\}\\
     remote = \{\}\quad unknown = []\quad uncovered = []};

  \foreach \a/\b in {l/p,p/i,i/w,w/s,s/m} \draw[fleche] (\a) -- (\b);
  \foreach \a/\b in {l/el,p/ep,i/ei,w/ew,s/es,m/em} \draw[lien] (\a) -- (\b);

  % Lecture du système de fichiers réel
  \node[douce,text width=24mm,font=\sffamily\scriptsize,minimum height=9mm] (fs) at (-3.9,-4*\D) {FS réel\\(lecture seule :\\\texttt{lstat}, \texttt{glob})};
  \draw[fleche,densely dashed] (fs) -- (s);
  \node[douce,text width=24mm,font=\sffamily\scriptsize,minimum height=9mm] (rt) at (-3.9,-2*\D) {Table de règles\\(données, chargée\\à la demande)};
  \draw[fleche,densely dashed] (rt) -- (i);

  % Les cinq rubriques de l'impact
  \node[decision,fill=dREV,text width=19mm] (c1) at (-2.3,-6.45*\D) {détruit\\[-2pt]{\tiny capturable ?}};
  \node[decision,fill=dALLOW,text width=19mm] (c2) at (0.2,-6.45*\D) {crée\\[-2pt]{\tiny rien à capturer}};
  \node[decision,fill=dASK,text width=19mm] (c3) at (2.7,-6.45*\D) {distant\\[-2pt]{\tiny surface + sorte}};
  \node[decision,fill=dCTX,text width=19mm] (c4) at (5.2,-6.45*\D) {non compris\\[-2pt]{\tiny unknown}};
  \node[decision,fill=gris,text width=19mm] (c5) at (7.7,-6.45*\D) {non couvert\\[-2pt]{\tiny uncovered}};
  \draw[fleche] (m.south) -- ++(0,-0.45) coordinate (b);
  \draw[encre!70,line width=0.7pt] (c1.north |- b) -- (c5.north |- b);
  \foreach \n in {c1,c2,c3,c4,c5} \draw[fleche] (\n.north |- b) -- (\n.north);
  \node[etiquette,below=1mm of c3] {$\downarrow$ \texttt{verdict.py} puis \texttt{decide.py} (figure~\ref{fig:decision})};
\end{tikzpicture}
}
\caption{L'intérieur du moteur, sur un exemple réel (sortie du moteur au commit \commitref{})}
\label{fig:moteur}
\end{figure}

\subsection{Les étapes}
\begin{description}[style=nextline,leftmargin=0.6cm,font=\sffamily\bfseries]
  \item[Analyse syntaxique (\fichier{model.py})]
    Un analyseur lexical écrit pour le shell POSIX produit des \cmd{Word} qui gardent leur
    texte source, leur valeur sans citations, les positions de leurs citations et de leurs
    substitutions, et un indicateur « s'étend ». Les substitutions \cmd{\$(\ldots)},
    \cmd{`\ldots`} et \cmd{<(\ldots)} sont analysées comme des \emph{scripts imbriqués}, les
    \textit{heredocs} sont rattachés à leur redirection, les enveloppes (\cmd{sudo},
    \cmd{env}, \cmd{xargs}, \cmd{timeout}, \cmd{nice}\ldots) sont traversées pour trouver le
    vrai mot de commande, et la charge d'un \cmd{bash -c} est analysée à son tour.
  \item[Invocations et grammaires (\fichier{invocation.py})]
    Chaque commande devient une invocation : un outil, des options canoniques, des
    opérandes, des redirections. Les options sont lues avec la grammaire propre à l'outil
    (options groupées, valeurs collées comme \cmd{sed -i.bak}, permutation à la GNU).
  \item[Expansion des mots (\fichier{words.py})]
    L'expansion se fait \emph{sur les mots}, comme bash, et non sur le texte : citations
    simples littérales, doubles citations sans découpage, \cmd{\$NOM} connu seulement s'il a
    été affecté plus tôt sur la ligne, accolades puis \textit{globs} sur le texte non cité,
    évalués sur l'état simulé. Tout le reste lève \cmd{Unknown} avec sa raison --- rien n'est
    deviné.
  \item[Simulation (\fichier{simulate.py})]
    Les effets sont joués dans l'ordre sur une \emph{surcouche} du système de fichiers réel,
    qui n'est que lu (\cmd{lstat}, \cmd{listdir}). Chaque nœud retient l'origine de son
    contenu. Les séparateurs sont modélisés comme des \emph{ensembles d'états possibles} :
    après \cmd{a \&\& b}, \cmd{b} part de l'état où \cmd{a} a réussi ; après \cmd{a || b}, de
    l'état où \cmd{a} a échoué --- avant \cmd{a} ou en partie ; les fusions sont des unions,
    donc elles ne peuvent qu'\emph{ajouter} à ce qui est détruit. Les sous-shells copient le
    répertoire courant et les variables mais gardent leurs effets sur les fichiers.
  \item[Impact (\fichier{effects.py})]
    Le résultat : \cmd{destroys} (indexé par le chemin \emph{préexistant} dont le contenu est
    perdu), \cmd{creates}, \cmd{touched} (tout chemin changé, pour le plan de contrôle),
    \cmd{remote} (surface $\rightarrow$ sorte d'effet), \cmd{unknown}, \cmd{uncovered}, et
    les bases de données capturables ou non.
\end{description}

\subsection{Le langage des effets}
Un vocabulaire unique décrit ce qu'une commande fait au monde (tableau~\ref{tab:effets}).
Quand une même surface distante est atteinte deux fois, elle garde la sorte qui perd le
plus, selon l'ordre DELETE $>$ OVERWRITE $>$ SEND $>$ RUN $>$ LATER $>$ SERVICE\_LOSS $>$
REF $>$ SERVICE\_KEEP $>$ CREATE.

\begin{table}[htbp]
\centering
\caption{Le langage des effets (\fichier{engine/effects.py})}
\label{tab:effets}
\small
\begin{tabularx}{\linewidth}{@{}L{2.8cm}YL{4.2cm}@{}}
\toprule
\textbf{Effet} & \textbf{Signification} & \textbf{Conséquence} \\
\midrule
DELETE, OVERWRITE, TRUNCATE & un chemin est supprimé, remplacé, vidé & copier avant \\
MOVE & un contenu passe de la source à la destination & copier la destination si elle existe \\
CREATE, APPEND, MKDIR & un chemin apparaît, ou reçoit un ajout & rien de préexistant n'est perdu \\
REF & une référence git bouge & le \textit{reflog} la ramène (arbitrage c) \\
SEND & part vers l'extérieur sans rappel possible & escalade \\
RUN & du code s'exécute, ici ou là-bas & lu à travers ce code, ou filet \\
LATER & du code est planté pour s'exécuter plus tard & escalade (ask) \\
SERVICE\_LOSS / KEEP & disponibilité perdue / gardée & escalade / autorisé et listé \\
UNKNOWN & quelque chose qu'on ne peut résoudre & la raison est gardée \\
\bottomrule
\end{tabularx}
\end{table}

\subsection{La table de règles, écrite comme des données}
Une règle (\fichier{engine/rules.py}) associe un ou plusieurs noms d'outil à une grammaire
d'options, une condition sur les options, des effets, et des \textbf{exemples}. Les exemples
ne sont pas de la documentation : chacun devient un test qui exécute tout le moteur dans une
arborescence fixe et compare l'impact obtenu à l'impact attendu. Un outil que les gabarits ne
peuvent pas exprimer (\cmd{git}, \cmd{aws}, \cmd{rsync}\ldots) délègue à un
\emph{résolveur}, mais garde dans la table sa grammaire et ses exemples.

\begin{code}{Extrait réel de la règle \cmd{rm} (\fichier{engine/rules.py})}
Rule(
    id="rm",
    tools=("rm",),
    options=OptionSpec(flags={
        "-f": "force", "--force": "force",
        "-r": "recursive", "-R": "recursive", "--recursive": "recursive",
        "-d": "dir", "--dir": "dir",
        # ...
    }),
    # effets et exemples : chaque exemple est un test
)
\end{code}

Au commit de référence, la table compte \nbregles{} règles, \nbexemples{} exemples et
\nboutils{} noms d'outils, répartis entre un noyau toujours chargé (fichiers locaux,
redirections) et \nbfamilles{} familles chargées \emph{à la demande} : SQL, git,
infrastructure, deux familles cloud, conteneurs, entrepôts de données, envois sortants,
PaaS, SaaS et migrations, système, coffres de secrets. Le chargement paresseux tient le
coût d'import hors du chemin de chaque commande : une ligne \cmd{rm} ne paie jamais
l'analyseur SQL. La même table génère aussi la liste des mots que le piège du shell envoie au
moteur (partie M2.2e) : une famille ajoutée atteint la porte du shell sans ligne de code
supplémentaire. L'annexe~\ref{ann:regles} reproduit une règle réelle de la table.

\subsection{Le budget}
Le moteur s'exécute sur le chemin critique de chaque commande de l'agent. Il dispose de
10\,ms de temps CPU (mesuré en temps processeur sous Linux, où l'horloge est fine), d'une
seconde tentative à 50\,ms en cas de dépassement, et d'un plafond de 250\,ms en temps réel.
Le temps d'import d'une famille et la lecture de certains fichiers sont exclus du budget.
S'il ne termine pas, la règle de la section~\ref{sec:menace} s'applique : jamais une
décision plus permissive.

% 3.6 — La politique et le décideur unique
\section{La politique et le décideur unique}
\label{sec:decision}

\subsection{Une seule politique}
Avant la partie T5, les mêmes questions --- « une suppression distante escalade-t-elle ? »,
« quelle surface derrière l'étiquette \cmd{db:} ? », « où restaurer un bucket ? » ---
avaient des réponses dans plusieurs modules. Elles sont désormais réunies dans
\fichier{policy.py}, que lisent le décideur, le jugement du moteur, l'orchestrateur et
l'affichage du modèle de menace :
\begin{itemize}
  \item \cmd{REMOTE\_POLICY} : sorte d'effet distant $\rightarrow$ ALLOW ou ESCALATE ; une
    sorte inconnue escalade ;
  \item \cmd{\_LABEL\_SURFACES} : étiquette du moteur $\rightarrow$ surface
    (\cmd{storage:} $\rightarrow$ \cmd{object\_storage}, \cmd{vcs-remote:} $\rightarrow$
    \cmd{remote\_vcs\_history}\ldots) ;
  \item \cmd{ESCALATION\_CHANNEL} : surface $\rightarrow$ ask ou deny ;
  \item \cmd{RECOVERY\_HINTS} : surface $\rightarrow$ où le fournisseur garde son propre
    historique (« bucket versioning, if enabled », « the database's PITR / snapshot »).
\end{itemize}

\subsection{La table de décision}
La figure~\ref{fig:decision} résume ce que la politique et le décideur produisent selon
\emph{où} l'effet tombe et \emph{quelle sorte} d'effet il est. Chaque case a été vérifiée en
faisant passer des commandes réelles par \cmd{Guard.evaluate}.

\begin{figure}[htbp]
\centering
% Figure 4 — La table de décision : où l'effet tombe x quelle sorte d'effet -> décision.
% Source : policy.py (REMOTE_POLICY, ESCALATION_CHANNEL, NO_SURFACE_CHANNEL),
% engine/verdict.py, decide.py, guard._unread_facts (M2.3a).
\begin{tikzpicture}[
    cell/.style={minimum width=30mm,minimum height=14mm,draw=white,line width=1pt,
                 fill=grisclair,inner sep=1.5pt,align=center,text width=28mm,font=\sffamily\tiny},
    chip/.style={decision,minimum height=4.6mm,inner sep=1.5pt,font=\sffamily\tiny\bfseries},
    note/.style={font=\sffamily\tiny,text=gris,align=center,text width=28mm},
    tete/.style={font=\sffamily\scriptsize\bfseries,align=center,text width=29mm},
    rang/.style={font=\sffamily\scriptsize,align=left,text width=31mm},
  ]
  \def\W{3.05}\def\H{1.5}
  \def\cellule#1#2#3#4{% x y chips note
    \node[cell] at (#1*\W,-#2*\H) {#3\\[1pt]{\color{gris}#4}};}
  \def\cA{\puce{dALLOW}{ALLOW}}
  \def\cR{\puce{dREV}{REVERSIBLE}}
  \def\cS{\puce{dSUP}{SUPERVISED}}
  \def\cEa{\puce{dASK}{ESCALATE\,·\,ask}}
  \def\cEd{\puce{dDENY}{ESCALATE\,·\,deny}}

  % En-têtes
  \node[tete] at (1*\W,0.15) {Ajout\\{\tiny\mdseries CREATE, APPEND, SERVICE\_KEEP}};
  \node[tete] at (2*\W,0.15) {Perte\\{\tiny\mdseries DELETE, OVERWRITE, TRUNCATE, MOVE, SERVICE\_LOSS}};
  \node[tete] at (3*\W,0.15) {Irrévocable\\{\tiny\mdseries SEND, RUN}};
  \node[tete] at (4*\W,0.15) {Plus tard\\{\tiny\mdseries LATER}};

  \node[rang] at (-0.05,-1*\H) {\textbf{Dans le projet}\\{\color{gris}fichiers, base sqlite}};
  \node[rang] at (-0.05,-2*\H) {\textbf{Hors du projet}\\{\color{gris}même machine}};
  \node[rang] at (-0.05,-3*\H) {\textbf{Fichiers de \texttt{demo\_cli}}\\{\color{gris}plan de contrôle}};
  \node[rang] at (-0.05,-4*\H) {\textbf{Surface « ask »}\\{\color{gris}base distante, git distant, bucket, conteneur\ldots}};
  \node[rang] at (-0.05,-5*\H) {\textbf{Surface « deny »}\\{\color{gris}infra détruite, flush, disque, paiement\ldots}};

  % Projet
  \cellule{1}{1}{\cA}{rien d'existant détruit}
  \cellule{2}{1}{\cR}{capture exacte, tout ou rien ; sinon ESCALATE·ask}
  \cellule{3}{1}{\cS}{code non lu : filet, sinon point de contrôle, sinon ALLOW + note (ESCALATE·ask pour code en ligne et installations, Rb2)}
  \cellule{4}{1}{\cEa}{hook git, \texttt{.envrc}\ldots}
  % Hors projet
  \cellule{1}{2}{\cA}{rien d'existant détruit}
  \cellule{2}{2}{\cEa}{aucune copie possible (V17)}
  \cellule{3}{2}{\cS}{le filet reste (V21) ; hors projet il enregistre sans sauvegarder (V18)}
  \cellule{4}{2}{\cEa}{\texttt{\textasciitilde/.bashrc}, crontab, tâche\ldots}
  % Plan de contrôle
  \cellule{1}{3}{\cEd}{même une création}
  \cellule{2}{3}{\cEd}{}
  \cellule{3}{3}{\cEd}{}
  \cellule{4}{3}{\cEd}{}
  % Distant ask
  \cellule{1}{4}{\cA}{listé dans le reçu}
  \cellule{2}{4}{\cEa}{indice : PITR, versions, reflog\ldots}
  \cellule{3}{4}{\cEa}{\texttt{ssh hôte cmd}, publication, création qui porte du code (Rb4)}
  \cellule{4}{4}{\cEa}{\texttt{deferred\_execution}}
  % Distant deny
  \cellule{1}{5}{\cA}{listé dans le reçu}
  \cellule{2}{5}{\cEd}{tout part, pas de retour réaliste}
  \cellule{3}{5}{\cEd}{paiement, file purgée\ldots}
  \cellule{4}{5}{}{---}

  % Lignes transverses
  \node[rang,text width=150mm,anchor=west] at (-1.6,-6.05*\H) {%
    \textbf{Ligne non comprise :} filet $\rightarrow$ point de contrôle $\rightarrow$ \puce{dASK}{ESCALATE\,·\,ask}
    \quad \textbf{Télécharger-et-exécuter :} \puce{dDENY}{ESCALATE\,·\,deny}\\[2pt]
    \textbf{Cible déclarée « production » :} capture, puis \puce{dCTX}{CONTEXT\_MISMATCH} montré à la personne, jamais approuvable sur place.};
\end{tikzpicture}

\caption{La table de décision : où l'effet tombe $\times$ quelle sorte d'effet}
\label{fig:decision}
\end{figure}

\subsection{Le décideur}
\fichier{decide.py} est une fonction pure : elle ne lit aucun fichier, ne capture rien, et
reçoit tout ce dont elle a besoin --- la classification, l'environnement, le fait qu'une
capture de la \emph{vraie} cible a réussi, les désaccords de contexte, la disponibilité d'un
filet. Ses branches s'évaluent dans un ordre fixe (tableau~\ref{tab:ordre}) ; la première qui
s'applique décide.

\begin{table}[htbp]
\centering
\caption{L'ordre des branches de \fichier{decide.py}}
\label{tab:ordre}
\small
\begin{tabularx}{\linewidth}{@{}L{0.6cm}YL{3.6cm}@{}}
\toprule
& \textbf{Situation} & \textbf{Décision} \\
\midrule
1a & Code illisible exécuté sur une surface nommée (\cmd{ssh hôte cmd}, \cmd{gh workflow run}) &
  ESCALATE, canal de la surface \\
1b & Télécharger et exécuter (\cmd{curl \ldots{} | bash}) & ESCALATE, deny \\
2 & Rien ne modifie & ALLOW \\
3 & Surface non récupérable & ESCALATE (ALLOW avec jeton structurel) \\
4 & Désaccord avec le contexte déclaré & CONTEXT\_MISMATCH si capturé, sinon ESCALATE \\
4b & Cible \emph{déclarée} production, capturée & CONTEXT\_MISMATCH \\
5 & Capture réussie & DRY\_RUN (aperçu SQL) ou REVERSIBLE \\
6 & Un filet couvre le projet & SUPERVISED \\
7 & Sinon & ESCALATE (ask) \\
\bottomrule
\end{tabularx}
\end{table}

\begin{pourquoi}[Pourquoi un seul décideur ?]
Avec deux chemins de décision, il y a deux réponses possibles à la même question ; c'est la
forme même des défauts trouvés par l'audit. Le moteur, l'ancien chemin, la porte fichier et
le chemin des lignes non lues ne font que \emph{rassembler des faits} ; seul
\fichier{decide.py} conclut. Une correction de politique s'applique donc partout à la fois,
sur les deux systèmes, et le portail de décision n'a qu'un seul point à observer.
\end{pourquoi}

\subsection{Six décisions, trois postures}
Six décisions sont précises mais techniques. Pour la personne, elles sont regroupées en trois
postures : \textbf{SAFE} (ALLOW, DRY\_RUN, REVERSIBLE, SUPERVISED --- la commande
s'exécute), \textbf{REVIEW} (CONTEXT\_MISMATCH --- capturé, puis montré), \textbf{BLOCKED}
(ESCALATE). La décision exacte reste affichée en sous-titre, comme dans les sorties de
terminal de ce chapitre (« \texttt{[+] SAFE · REVERSIBLE} »).

\subsection{Les escalades sans surface}
Trois escalades ne portent pas de surface et ont leur propre canal
(\cmd{NO\_SURFACE\_CHANNEL}) : une modification locale que rien n'a pu capturer (ask), un
désaccord de contexte sans capture (ask), et le téléchargement-exécution (deny), dont
personne ne peut voir la charge avant qu'elle ne s'exécute.

\begin{terminal}[une surface distante : la personne est consultée, l'indice de restauration est donné]
$ demo_cli check --mode enforce 'git push --force origin main'
demo_cli 1.7.0  [x] BLOCKED  ·  ESCALATE
  matched rule        git_force_push
  surface             remote_vcs_history
Decision
  Mutation of a non-recoverable surface (remote_vcs_history); a local snapshot
  cannot cover its blast radius.
Safe next step
  1. Confirm the external effect is intended and authorised.
  2. Supply a structural approval token, or perform it outside the agent.
  3. If it already ran, restore from the provider: a teammate's
     un-force-pushed clone or the remote's reflog.
\end{terminal}

% 3.7 — Deux systèmes, un décideur
\section{Deux systèmes, un décideur}
\label{sec:deux-os}

Les agents ne travaillent pas de la même façon sous Linux et sous Windows : sous Windows,
Claude Code dispose d'un outil \cmd{PowerShell} à côté de son outil \cmd{Bash} (Git Bash), et
Codex écrit du PowerShell. Les chemins s'écrivent \cmd{C:\textbackslash x} ou \cmd{/c/x}, il
n'y a pas de ptrace, et les appels système n'ont rien de commun. La conception sépare donc
ce qui \emph{doit} être commun de ce qui \emph{ne peut pas} l'être
(figure~\ref{fig:os}).

\begin{figure}[htbp]
\centering
% Figure 5 — Linux et Windows : un seul décideur, deux filets d'exécution.
\begin{tikzpicture}[
    col/.style={boite,text width=57mm,font=\sffamily\footnotesize},
    titre/.style={font=\sffamily\bfseries\small,text=encre},
  ]
  % Partie commune
  \node[accent,text width=72mm,minimum height=10mm] (eng) at (0,0)
     {Moteur partagé : simulation, effets, impact\sous{même table de règles, même \texttt{policy.py}}};
  \node[boite,draw=encre,line width=1.1pt,text width=72mm,below=4mm of eng] (dec)
     {\textbf{\texttt{decide.py}} --- un seul décideur\sous{mêmes garanties sur les deux systèmes}};
  \draw[fleche] (eng) -- (dec);

  % Frontaux
  \node[boite,text width=40mm] (fp) at (-4.75,1.9) {Frontal POSIX\sous{\texttt{model.py} ; Git Bash : \texttt{/c/x} $\rightarrow$ \texttt{C:\textbackslash x}}};
  \node[boite,text width=40mm,dashed] (fw) at (4.75,1.9) {Frontal PowerShell (C3)\sous{lieur de paramètres ; \statutCthree}};
  \draw[fleche] (fp.south) |- (eng.west);
  \draw[fleche] (fw.south) |- (eng.east);

  % Colonnes OS
  \node[titre] (tl) at (-4.75,-2.9) {Linux};
  \node[titre] (tw) at (4.75,-2.9) {Windows};
  \node[col,below=2mm of tl] (l1) {Agents : Claude Code (Bash), Codex,\\piège bash (\texttt{BASH\_ENV}, \texttt{trap DEBUG})};
  \node[col,below=2mm of l1,fill=wepwnclair,draw=wepwnfonce] (l2) {\textbf{Garde d'appels système}\sous{\texttt{ptrace} : arrêt à l'entrée de \texttt{unlink}, \texttt{rename},\\\texttt{open(O\_TRUNC)}\ldots{} ; copie \emph{puis} laisse faire}};
  \node[col,below=2mm of l2] (l3) {Activation : réécriture par le \textit{hook}\\\texttt{demo\_cli sh -g <id> <cmd>} (sans \texttt{allow})};
  \node[col,below=2mm of l3] (l4) {Plan de contrôle : appel annulé\\(\texttt{orig\_rax = -1}) + reçu ESCALATE};

  \node[col,below=2mm of tw] (w1) {Agents : Claude Code (Bash \& PowerShell),\\Codex (PowerShell)};
  \node[col,below=2mm of w1,fill=wepwnclair,draw=wepwnfonce] (w2) {\textbf{Montage WinFsp}\sous{système de fichiers en espace utilisateur : \texttt{rename}\\(remplacement), \texttt{cleanup}(DELETE), \texttt{set\_file\_size}}};
  \node[col,below=2mm of w2] (w3) {Activation : le projet \emph{est} le montage\\(\texttt{C:\textbackslash p} $\rightarrow$ données dans \texttt{C:\textbackslash p.real} verrouillé)};
  \node[col,below=2mm of w3] (w4) {Remontage à l'ouverture de session :\\tâche élevée durcie (W-SEC)};

  \draw[fleche] (dec.south) -- ++(0,-0.5) -| (l1.north);
  \draw[fleche] (dec.south) -- ++(0,-0.5) -| (w1.north);

  \node[zone,fit=(tl)(l1)(l4)] (zl) {};
  \node[zone,fit=(tw)(w1)(w4)] (zw) {};
  \node[etiquette,text width=120mm,anchor=north] at ([yshift=-2mm]0,0 |- zw.south) {Les deux filets couvrent le \emph{projet} seulement, et la raison de la décision le dit.};
\end{tikzpicture}

\caption{Linux et Windows : un seul décideur, deux filets d'exécution}
\label{fig:os}
\end{figure}

\begin{description}[style=nextline,leftmargin=0.6cm,font=\sffamily\bfseries]
  \item[Commun]
    Le moteur, la table de règles, la politique, le décideur, le format des reçus, la
    récupération et \cmd{undo}. Une ligne POSIX tapée dans Git Bash sous Windows passe par le
    même moteur que sous Linux ; ses chemins \cmd{/c/\ldots} sont traduits en
    \cmd{C:\textbackslash\ldots} (partie M2.1), et un chemin POSIX absolu dont la racine MSYS
    n'est pas connue (\cmd{/tmp}, \cmd{/usr}) est marqué non compris plutôt que deviné.
  \item[Propre à chaque système]
    Le filet d'exécution et son activation, la protection du plan de contrôle pendant
    l'exécution, et le frontal de lecture : POSIX d'un côté, PowerShell de l'autre (C3,
    section~\ref{sec:c3}).
\end{description}

\subsection{Linux : pourquoi ptrace}
Le filet Linux est un \emph{garde d'appels système} fondé sur ptrace
(\fichier{syscall\_guard.py}). Il lance la commande comme processus suivi et s'arrête à
l'\emph{entrée} de chaque appel système ; pour un appel destructeur (\cmd{unlink},
\cmd{unlinkat}, \cmd{rmdir}, \cmd{rename*} sur une destination existante, \cmd{truncate},
\cmd{open} avec \cmd{O\_TRUNC} ou en écriture) sur un fichier préexistant du projet, il copie
le fichier \emph{puis} laisse l'appel se faire. Pour un fichier du plan de contrôle, il annule
l'appel.

Trois mécanismes étaient envisageables :
\begin{itemize}
  \item \textbf{FUSE} : monter un système de fichiers sous le projet. Il a été écarté parce que
    le garde d'appels système fait déjà la même chose --- voir chaque destruction dans le
    projet avant qu'elle n'ait lieu. De plus, un montage FUSE suppose de placer le projet
    derrière un point de montage servi par un démon permanent, alors que ptrace ne s'attache
    qu'à la commande concernée.
  \item \textbf{seccomp avec notification} : plus rapide, car seuls les appels filtrés
    réveillent le superviseur. Le code le désigne comme la cible de production du filet ; le
    prototype actuel utilise ptrace, qui donne la main à l'entrée de chaque appel et permet de
    lire les registres et la mémoire du processus suivi. \acompleter[confirmer la raison du choix de ptrace pour le prototype]
  \item \textbf{ptrace} : retenu. Ses limites sont assumées et rappelées au
    chapitre~\ref{ch:securite} : architecture x86-64, coût d'un arrêt par appel système, et
    pas de vision au-delà du projet.
\end{itemize}

\paragraph{Activation.} Le filet doit s'exécuter \emph{à la place} de la commande. Le
\textit{hook} Claude Code réécrit donc la commande en \cmd{demo\_cli sh -g <id> <commande>}
via \cmd{updatedInput} --- \emph{sans} répondre \cmd{allow}, pour que l'invite de permission
de l'utilisateur continue de s'afficher (partie B1b) --- et dit à la personne et au modèle ce
qui s'exécute réellement. Un \cmd{demo\_cli sh} tapé par l'agent lui-même est déballé et jugé
sur ce qu'il cache ; un processus déjà suivi par un traceur étranger refuse de s'exécuter
sans surveillance.

\subsection{Windows : pourquoi un montage WinFsp}
Windows n'a pas de ptrace. L'endroit où toutes les routes vers un fichier convergent --- API
Win32, \cmd{ntdll}, appels système directs --- est le système de fichiers. \demo{} y place un
système de fichiers en espace utilisateur WinFsp (\fichier{fsmount.py}) : le projet
\cmd{C:\textbackslash projet} devient le point de montage, et les données vivent dans
\cmd{C:\textbackslash projet.real}, verrouillé par une ACL réservée aux Administrateurs et à
SYSTEM. Le montage s'exécute avec élévation, l'agent sans : l'agent ne peut pas contourner le
montage en écrivant directement dans le répertoire réel (\fichier{protect.py}). Le filet
n'exige donc aucune réécriture de commande : tout ce qui touche le projet passe par lui.

La logique de jugement (\fichier{fsguard.py}) est séparée de la plomberie WinFsp et testable
sous Linux. Ses règles viennent de journaux d'opérations \emph{observés} sous Windows~10,
qui ont contredit l'intuition sur trois points sur quatre : le remplacement atomique
(\cmd{rename} avec remplacement) détruit la destination sans aucun signal de suppression ---
c'est ainsi que l'outil \cmd{Write} de Claude Code enregistre chaque fichier ; la méthode
\cmd{overwrite()} n'est jamais appelée, les vraies réécritures passant par une réduction de
taille ; et seul \cmd{cleanup} avec le bit de suppression signale toutes les suppressions.

Un montage meurt avec la session. Une tâche planifiée à l'ouverture de session le relance ;
elle s'exécute avec élévation, ce qui en fait une cible : sa conception durcie (W-SEC) est
décrite à la section~\ref{sec:wsec}.

\subsection{Points de contrôle : utiles, mais optionnels}
Quand ni capture ciblée ni filet ne sont possibles, \demo{} peut copier tout le projet avant
la commande (\fichier{checkpoint.py}) --- « je ne sais pas ce que vous allez détruire, donc je
préserve tout ce que vous pourriez détruire ». Cette option est \textbf{désactivée par
défaut} : sur un dépôt avec un gros \fichier{.git}, la copie est lente, et « un garde-fou
qui devient lent sans qu'on le lui demande est un garde-fou qu'on désinstalle ». Elle est
bornée (256\,Mo et \num{25000} fichiers par défaut), refuse le répertoire personnel et la
racine, et ne remplace jamais une capture ciblée ni ne couvre une surface distante ou un
fichier hors du projet.

% 3.8 — Choix de conception et alternatives
\section{Choix de conception et alternatives écartées}
\label{sec:choix}

Le tableau~\ref{tab:choix} rassemble les décisions d'architecture, l'alternative écartée et
la raison. Les justifications détaillées se trouvent dans les sections citées.

{\small
\begin{longtable}{@{}L{3.6cm}L{3.6cm}L{7.2cm}@{}}
\caption{Choix de conception, alternatives et justifications}\label{tab:choix}\\
\toprule
\textbf{Choix} & \textbf{Alternative écartée} & \textbf{Pourquoi} \\
\midrule
\endfirsthead
\toprule
\textbf{Choix} & \textbf{Alternative écartée} & \textbf{Pourquoi} \\
\midrule
\endhead
\bottomrule
\endfoot
Analyse syntaxique + simulation & Expressions régulières &
  Le shell est une grammaire ; les motifs se trompent dans les deux sens. L'analyse unique a
  corrigé 22 cas sans en casser aucun (§\ref{sec:pourquoi-moteur}). \\
Une seule analyse lue par tous & Un analyseur par besoin &
  Deux analyseurs qui répondent différemment à la même question étaient la source principale
  des défauts de l'audit. \\
Logique déterministe sur la commande & Un LLM qui juge la commande &
  L'outil accompagne des agents fondés sur des LLM ; il doit donner la même réponse à la même
  ligne, être testable exemple par exemple, tenir en quelques millisecondes et ne pas
  pouvoir être convaincu par le texte qu'il analyse. \\
Règles écrites comme des données, exemples = tests & Code ad hoc par outil &
  Une règle se lit, se relit et se vérifie ; ses \nbexemples{} exemples sont exécutés à chaque
  test ; le préfiltre du shell en est généré. \\
Grammaires tirées des outils eux-mêmes & Grammaires écrites de mémoire &
  Une option mal connue est une décision fausse (\cmd{unzip -Pxo}, \cmd{helm 4},
  \cmd{tofu -chdir}) ; les règles non vérifiées sont marquées (§\ref{sec:faits-verifies}). \\
Récupérer d'abord & Bloquer tout ce qui est dangereux &
  Héritage de la bêta : bloquer arrête l'agent et pousse à affaiblir les règles ; une copie
  ciblée laisse le travail continuer et rend l'erreur réversible. \\
Capture ciblée, tout ou rien & Copie partielle « au mieux » &
  REVERSIBLE est une promesse ; une copie partielle la rend fausse (§\ref{sec:jamais-reversible}). \\
Points de contrôle optionnels & Copie du projet avant chaque commande &
  Coût prohibitif sur un gros \fichier{.git} ; un outil lent est désinstallé. \\
Moteur en ombre, puis décideur & Basculer directement &
  Les divergences entre prédiction et décision ont été comprises une à une, à décisions
  constantes, avant que le moteur ne décide (S2+M1 : 124 changements revus). \\
L'ancien chemin garde le plus strict, puis le moteur lit tout & Supprimer les anciennes règles d'un coup &
  Aucune décision ne doit devenir plus permissive pendant la migration ; M2.3b a terminé la sortie de l'ancien chemin pour POSIX, validée sur un corpus élargi de \nbcorpuslarge{} commandes (§\ref{sec:plus-strict}). \\
Un décideur pur et unique & Décisions dans chaque porte &
  Une correction s'applique partout ; le décideur ne peut pas revendiquer une capture qu'il
  n'a pas reçue. \\
Canal ask/deny par surface, deny par défaut & Tout demander, ou tout refuser &
  L'attention s'use ; une surface nouvelle ne devient jamais « demandable » par accident
  (§\ref{sec:ask-deny}). \\
Filet Linux par ptrace & FUSE &
  Le garde d'appels système fait déjà ce que ferait un montage FUSE, sans placer le projet
  derrière un démon permanent (§\ref{sec:deux-os}). \\
Filet Windows par WinFsp & Pilote de filtre noyau &
  WinFsp donne un système de fichiers en mode utilisateur, sans pilote noyau à écrire ni à
  signer ; il voit toutes les routes vers les fichiers. \\
Tâche planifiée à l'ouverture de session & Service Windows &
  Une tâche se crée et se supprime en une commande et reste visible dans le Planificateur ;
  rien de ce que l'outil fait à la machine n'est caché. Elle est durcie (W-SEC). \\
C3 : analyseur PowerShell maison & Appeler \cmd{powershell.exe} pour analyser chaque ligne &
  Lancer PowerShell à chaque commande coûte des centaines de millisecondes et le garde doit
  aussi fonctionner depuis WSL ; l'analyseur de PowerShell sert d'\emph{oracle} dans les
  tests (§\ref{sec:c3}). \\
C3 : tables de paramètres générées & Listes écrites à la main &
  PowerShell a des centaines de paramètres, alias et abréviations ; une nouvelle version
  devient une régénération, pas une modification de code. \\
\end{longtable}
}


\section*{Conclusion du chapitre}
L'architecture tient en une phrase : \emph{une seule lecture, une seule politique, un seul
décideur, et des filets là où la lecture s'arrête}. Les principes qui la gouvernent
partagent une même exigence --- ne jamais dire plus que ce qui est vrai : pas d'ALLOW pour
une ligne non lue, pas de REVERSIBLE pour une capture incomplète, pas de « demander » là où
personne ne verra la question. Le chapitre suivant montre comment ces décisions sont
devenues du code.
